\documentclass[UTF8,11pt,a4paper]{ctexart}
\usepackage[a4paper,margin=25mm]{geometry}
\usepackage{amsmath,amssymb,amsthm,mathtools,bm}
\usepackage{booktabs,longtable,array}
\usepackage{enumitem}
\usepackage{microtype}
\usepackage{hyperref}
\usepackage{xcolor}
\usepackage{fancyhdr}
\usepackage{setspace}

\definecolor{深蓝}{RGB}{33,72,108}
\hypersetup{colorlinks=true,linkcolor=深蓝,urlcolor=深蓝}
\pagestyle{fancy}
\fancyhf{}
\fancyhead[L]{\small 催化相容性交互图册}
\fancyhead[R]{\small 数学建模与理论基础}
\fancyfoot[C]{\thepage}
\setlength{\headheight}{14pt}
\setstretch{1.16}
\setlist{itemsep=0.25em,topsep=0.4em}

\newtheorem{definition}{定义}[section]
\newtheorem{theorem}[definition]{定理}
\newtheorem{proposition}[definition]{命题}
\newtheorem{corollary}[definition]{推论}
\theoremstyle{remark}
\newtheorem{remark}[definition]{说明}

\newcommand{\Rset}{\mathcal R}
\newcommand{\Eset}{\mathcal E}
\newcommand{\Cset}{\mathcal C}
\newcommand{\Kstar}{K^{\star}}
\newcommand{\norm}[1]{\left\lVert #1\right\rVert}
\newcommand{\abs}[1]{\left\lvert #1\right\rvert}

\title{\vspace{-1.2cm}\Huge\bfseries 催化酶—反应交互图册\\[0.55em]
\Large 数学建模、设计依据与核心证明}
\author{南京大学二〇二六年国际遗传工程机器大赛团队}
\date{二〇二六年九月二十七日}

\begin{document}
\maketitle

\begin{abstract}
酶检索面对一种特殊的数据结构：单个反应和蛋白拥有丰富的分子信息，可靠酶—反应配对却十分稀疏，许多局部信息还只在部分对象上可用。本文给出当前模型的完整数学表述。模型用若干局部交互图描述不同信息区域，每个局部图拥有自己的反应坐标、酶坐标和双线性交互形式；各图通过非负且和为一的权重组合。缺失信息对应的局部图权重严格为零。当前两个检索方向共享局部配对分数，同时保留不同方向的读出权重。本文证明坐标变换不改变局部交互、加权组合的凸包性质、单图精确退回、方向读出差异上界、重叠一致性损失的零点条件，以及有限秩实验更新的秩界和扰动界。
\end{abstract}

\tableofcontents

\section{建模对象从哪里来}

\subsection{真实催化活性带有条件}

令 $r\in\Rset$ 表示反应，$e\in\Eset$ 表示酶，$c\in\Cset$ 表示温度、酸碱度、辅因子、浓度、宿主、检测方式等实验或细胞条件。概念上，真实活性应写成
\begin{equation}
A:\Rset\times\Eset\times\Cset\rightarrow\mathbb R,
\qquad A(r,e;c).
\end{equation}

现有关系数据只覆盖很少的配对，而且许多记录缺少完整条件。模型因此选择一个较弱、可以由现有数据支持的目标。

\begin{definition}[工作催化相容性]
在给定候选范围、数据快照和观测语义下，记
\begin{equation}
\Kstar:\Rset\times\Eset\rightarrow\mathbb R
\end{equation}
为任务层的工作催化相容性。它用于排序与检索，不承诺覆盖所有实验条件下的真实活性。
\end{definition}

这一定义直接约束了负例语义。某个配对没有记录，通常只能说明“未知”；某个条件下检测不到活性，也只约束该条件。

\subsection{数据结构决定局部建模}

反应侧可以观测整体反应、底物和产物、反应指纹、原子映射、键变化、反应中心等信息；酶侧可以观测序列、蛋白语言模型表示、家族、结构、口袋、基序和辅因子等信息。不同信息的覆盖率差异很大，维度和语义也不同。

如果强制把这些信息先压入一个全局反应空间和一个全局酶空间，局部信息会遇到两个问题：覆盖稀疏时全局距离缺少实际观测基础；某些局部信息对一个方向有益，对另一个方向会破坏原有排序。历史上的反应中心切空间和口袋覆盖审计都出现过这种现象。因此当前模型直接在酶—反应配对上建立多个局部交互描述。

\section{局部交互图}

\begin{definition}[局部交互图]
第 $\alpha$ 个局部图由四部分组成：
\begin{equation}
U_\alpha\subseteq\Rset\times\Eset,
\end{equation}
表示它的适用区域；
\begin{equation}
\phi_\alpha(r)\in H_\alpha^R,
\qquad
\psi_\alpha(e)\in H_\alpha^E,
\end{equation}
分别表示反应坐标和酶坐标；矩阵 $G_\alpha$ 把两侧坐标配对；局部分数定义为
\begin{equation}
K_\alpha(r,e)=
\phi_\alpha(r)^{\mathsf T}G_\alpha\psi_\alpha(e).
\label{eq:local-score}
\end{equation}
\end{definition}

$H_\alpha^R$ 与 $H_\alpha^E$ 可以具有不同维度。两个不同局部图之间也不需要共享坐标轴。整体反应与序列可以形成一个图，反应中心与结构口袋可以形成另一个图，家族专项信息还可以形成第三个图。

\subsection{为什么选择双线性形式}

双线性形式把两侧分子坐标和“如何相互匹配”分开。坐标负责描述对象，矩阵 $G_\alpha$ 负责描述该局部区域里的催化需求与酶能力如何配对。它还允许两侧维度不同，适合异质分子表示。

若某个局部交互核具有快速衰减的奇异值谱，则其前若干项给出有限秩近似
\begin{equation}
K_\alpha(r,e)\approx
\sum_{j=1}^{d_\alpha}
\sigma_{\alpha j}u_{\alpha j}(r)v_{\alpha j}(e),
\end{equation}
这正对应双塔或低维局部专家的有限秩表达。这里的低秩是建模假设，需要继续接受实验检验。

\subsection{局部分数与坐标选择无关}

\begin{theorem}[基变换不变性]
设某个局部图采用列向量坐标，局部分数为
\begin{equation}
K=\phi^{\mathsf T}G\psi.
\end{equation}
对任意可逆矩阵 $P,Q$，令
\begin{equation}
\widetilde\phi=P\phi,
\qquad
\widetilde\psi=Q\psi,
\qquad
\widetilde G=P^{-\mathsf T}GQ^{-1}.
\end{equation}
则
\begin{equation}
\widetilde\phi^{\mathsf T}\widetilde G\widetilde\psi
=\phi^{\mathsf T}G\psi.
\end{equation}
\end{theorem}

\begin{proof}
直接代入：
\begin{align}
\widetilde\phi^{\mathsf T}\widetilde G\widetilde\psi
&=(P\phi)^{\mathsf T}P^{-\mathsf T}GQ^{-1}(Q\psi)\\
&=\phi^{\mathsf T}G\psi.
\end{align}
\end{proof}

因此局部图的科学对象是“坐标加交互形式”的组合。坐标轴本身不具有绝对意义。

\section{可用性、适用性与权重分解}

对配对 $(r,e)$，记 $a_\alpha(r,e)\in\{0,1\}$ 表示第 $\alpha$ 个局部图所需信息是否存在。再用 $q_\alpha(r,e)$ 表示该图的适用程度。可用局部图的权重可写成带掩码的归一化形式
\begin{equation}
\rho_\alpha(r,e)=
\frac{a_\alpha(r,e)e^{q_\alpha(r,e)}}
{\sum_\beta a_\beta(r,e)e^{q_\beta(r,e)}}.
\label{eq:partition}
\end{equation}
于是
\begin{equation}
\rho_\alpha\ge0,
\qquad
\sum_\alpha\rho_\alpha=1,
\end{equation}
而且 $a_\alpha=0$ 时必有 $\rho_\alpha=0$。

当前体系始终保留一个只依赖基础分子输入的通用图，因此任意有效反应和蛋白至少有一个可用局部图。

理想的配对级读出写成
\begin{equation}
\widehat K(r,e)=
\sum_\alpha\rho_\alpha(r,e)K_\alpha(r,e).
\label{eq:glued-score}
\end{equation}

\begin{theorem}[凸包性质]
对任意固定配对，式 \eqref{eq:glued-score} 满足
\begin{equation}
\min_{\alpha:\rho_\alpha>0}K_\alpha
\le \widehat K\le
\max_{\alpha:\rho_\alpha>0}K_\alpha.
\end{equation}
\end{theorem}

\begin{proof}
设活跃局部分数的最小值和最大值分别为 $m,M$。由 $m\le K_\alpha\le M$、$\rho_\alpha\ge0$ 以及 $\sum_\alpha\rho_\alpha=1$，有
\begin{equation}
m\le\sum_\alpha\rho_\alpha K_\alpha\le M.
\end{equation}
\end{proof}

\begin{corollary}[单图精确退回]
若某个配对只有一个局部图可用，则该图权重为 $1$，最终分数精确等于它的局部分数。
\end{corollary}

这个性质给缺失信息一个明确语义：局部图缺失时退出权重分解，其余图重新归一化。缺失本身不会产生负分。

\section{历史多专家形式如何落入当前图册}

历史实现使用一个通用通道和若干局部专家。若某个方向的专家总质量为 $m$，局部专家内部权重为 $g_k$，满足
\begin{equation}
g_k\ge0,\qquad \sum_k g_k=1,\qquad 0\le m\le1,
\end{equation}
历史分数为
\begin{equation}
S=(1-m)K_0+m\sum_k g_kK_k.
\end{equation}
定义
\begin{equation}
\rho_0=1-m,
\qquad
\rho_k=mg_k,
\end{equation}
便有
\begin{equation}
\rho_0+\sum_k\rho_k=1,
\qquad
S=\rho_0K_0+\sum_k\rho_kK_k.
\end{equation}

因此历史凸组合正好是单位权重分解的一个特例。这个重写本身不改变原有数值。

\section{当前实现中的方向化读出}

当前冻结实现共享所有局部分数 $K_\alpha(r,e)$，两个检索方向采用不同权重。记方向 $d$ 为反应找酶或酶找反应，则
\begin{equation}
S_d(r,e)=
\sum_\alpha\rho_\alpha^{(d)}(r,e)K_\alpha(r,e).
\label{eq:directional}
\end{equation}

反应找酶主要使用反应侧门控，酶找反应主要使用蛋白侧门控，两个方向还拥有独立的局部专家总质量。因此当前模型只要求两个方向估计同一个任务层相容性，不要求它们的最终数值逐点一致。

\begin{theorem}[方向读出差异上界]
设 $p,q$ 是同一组活跃局部分数 $K_1,\ldots,K_m$ 上的两个合法权重分布，
\begin{equation}
S_p=\sum_i p_iK_i,
\qquad
S_q=\sum_i q_iK_i.
\end{equation}
令 $K_{\max}$、$K_{\min}$ 分别为活跃局部分数的最大值和最小值，则
\begin{equation}
\abs{S_p-S_q}
\le
\frac12\norm{p-q}_1
\left(K_{\max}-K_{\min}\right).
\label{eq:direction-bound}
\end{equation}
\end{theorem}

\begin{proof}
因为 $\sum_i(p_i-q_i)=0$，可减去任意常数。取
\begin{equation}
c=\frac{K_{\max}+K_{\min}}2,
\end{equation}
则
\begin{align}
\abs{S_p-S_q}
&=\abs{\sum_i(p_i-q_i)(K_i-c)}\\
&\le\sum_i\abs{p_i-q_i}\abs{K_i-c}\\
&\le\frac{K_{\max}-K_{\min}}2\sum_i\abs{p_i-q_i}.
\end{align}
\end{proof}

式 \eqref{eq:direction-bound} 说明方向差异由两部分共同控制：权重差异和局部图之间的分数跨度。若局部图在重叠区域逐渐一致，方向化权重造成的最终差异也会受到压缩。

\section{重叠一致性损失}

局部图允许使用不同坐标，多个图同时适用时仍应尽量给出一致的任务层相容性。当前实现对每个样本计算归一化的加权平方差。

设活跃图集合为 $A$。对 $\alpha<\beta$ 定义
\begin{equation}
w_{\alpha\beta}
=\rho_\alpha\rho_\beta
\mathbf 1\{\alpha,\beta\text{均可用}\}.
\end{equation}
若总重叠权重
\begin{equation}
Z=\sum_{\alpha<\beta}w_{\alpha\beta}
\end{equation}
大于零，则单个样本的一致性损失为
\begin{equation}
L_{\mathrm{合}}=
\frac{\sum_{\alpha<\beta}
 w_{\alpha\beta}(K_\alpha-K_\beta)^2}
{Z}.
\label{eq:glue-loss}
\end{equation}
若 $Z=0$，该样本不进入一致性均值。训练时再对所有 $Z>0$ 的样本求平均。

归一化使损失关注“活跃重叠上的平均分歧”，减少局部图数量和总重叠质量对损失尺度的机械影响。

\begin{theorem}[一致性损失的零点]
对某个满足 $Z>0$ 的样本，构造重叠图：每个活跃局部图是一个节点，当 $w_{\alpha\beta}>0$ 时连接一条边。则 $L_{\mathrm{合}}=0$ 当且仅当每条边两端的局部分数相等。于是每个连通分量中的所有局部分数都相等。
\end{theorem}

\begin{proof}
式 \eqref{eq:glue-loss} 的分母为正，分子是若干非负项之和。分子为零等价于所有正权边对应的平方差均为零。沿连通路径传递即可得到同一连通分量内的局部分数相等。
\end{proof}

当前冻结训练对两个方向分别计算这项损失，再按主检索损失中的方向权重组合。一致性项使用较弱固定系数，以保持它对原检索目标的约束性而不过度主导训练。

\section{稀疏配对监督怎样进入模型}

单个对象的分子坐标主要来自序列、反应结构、分子描述和结构等对象侧信息。稀疏已验证配对承担四类作用：学习局部交互形式，学习局部图适用程度，约束重叠一致性，训练最终检索排序。

这一区分解释了双冷启动为何可行。只要一个新反应和一个新酶都能从原始分子输入得到通用图坐标，就能够计算
\begin{equation}
K_0(r^{\ast},e^{\ast})
\end{equation}
并获得有效读出。模型可打分性由分子输入及其坐标映射保证，训练表实体编号只承担数据索引作用。实际准确性仍需严格双冷启动和时间迁移评测验证。

\section{实验反馈的有界有限秩更新}

已确认的新阳性配对可以直接更新某个局部图的交互矩阵。设观测为
\begin{equation}
\Omega=\{(r_i,e_i,w_i)\}_{i=1}^{n},
\end{equation}
其中 $w_i\ge0$ 表示证据权重。把局部图权重并入系数
\begin{equation}
\eta_i=w_i\rho_\alpha(r_i,e_i),
\end{equation}
定义原始更新
\begin{equation}
\Delta_\alpha^{\mathrm{原}}
=\sum_{i=1}^{n}
\eta_i\phi_\alpha(r_i)\psi_\alpha(e_i)^{\mathsf T}.
\label{eq:update}
\end{equation}

实现会把整个矩阵统一缩放，使
\begin{equation}
\norm{\Delta_\alpha}_F\le\varepsilon_\alpha.
\end{equation}
更新后的局部交互矩阵为
\begin{equation}
G_\alpha^{\mathrm{新}}=G_\alpha+\Delta_\alpha.
\end{equation}

\begin{theorem}[秩界]
式 \eqref{eq:update} 的更新满足
\begin{equation}
\operatorname{rank}(\Delta_\alpha)\le n.
\end{equation}
若只有 $m$ 个系数 $\eta_i$ 非零，则更强地有
\begin{equation}
\operatorname{rank}(\Delta_\alpha)\le m.
\end{equation}
\end{theorem}

\begin{proof}
每个外积 $\phi_i\psi_i^{\mathsf T}$ 的秩至多为一，矩阵秩对加法满足次可加性。统一缩放不改变非零矩阵的秩。
\end{proof}

\begin{theorem}[单配对扰动界]
若
\begin{equation}
\norm{\phi_\alpha(r)}_2\le1,
\qquad
\norm{\psi_\alpha(e)}_2\le1,
\end{equation}
且 $\norm{\Delta_\alpha}_F\le\varepsilon_\alpha$，则更新造成的局部分数变化满足
\begin{equation}
\abs{
\phi_\alpha(r)^{\mathsf T}
\Delta_\alpha
\psi_\alpha(e)
}
\le\varepsilon_\alpha.
\end{equation}
\end{theorem}

\begin{proof}
由柯西不等式和谱范数不超过弗罗贝尼乌斯范数，
\begin{align}
\abs{\phi^{\mathsf T}\Delta\psi}
&\le \norm{\phi}_2\norm{\Delta}_2\norm{\psi}_2\\
&\le \norm{\Delta}_F\\
&\le \varepsilon.
\end{align}
\end{proof}

这两个定理给新实验的影响提供了明确控制。当前数学内核已经实现同维共享空间和两侧维度不同的矩形更新。实验记录解析、条件来源保存、适用图判定和跨图投影仍由上层流程承担。

\section{当前冻结实现与数学模型的对应}

当前冻结配方包含一个通用图和八个局部图。通用图负责基础分子输入的全覆盖，局部图承担不同表示或生化区域中的专项能力。两个方向共享通用和局部配对分数，各自拥有方向化权重。训练包含双向检索损失、门控均衡、熵、局部图多样性和弱重叠一致性约束。

这个实现与理想配对级图册之间存在一处清楚的差别：理想形式可以让权重直接依赖完整配对 $(r,e)$；当前冻结实现主要由查询侧产生权重。后续若研究配对级权重，需要重新冻结协议并验证两个方向的排序表现，数学对称性本身不足以支持替换。

\section{为什么最终选择这种几何}

当前图册对应五条长期实验事实。

第一，通用输入必须保证任何有效反应和蛋白都能进入检索，因此始终保留全覆盖通用图。

第二，结构、口袋、反应中心和家族信息的覆盖率与收益具有明显局部性，因此它们以局部图进入。

第三，多个局部信息源可以同时有效，也可能给出相互矛盾的判断，因此需要重叠一致性。

第四，两个检索方向的风险结构不同，历史实验多次证明自由强行对称会损害其中一方，因此当前保留方向化读出并用可证明上界约束其差异来源。

第五，新实验应该能够直接影响模型，同时需要限制局部证据对既有排序的扰动范围，因此采用有界有限秩更新。

\section{结论与边界}

当前模型的数学核心可以压缩为四个式子：
\begin{align}
K_\alpha(r,e)
&=\phi_\alpha(r)^{\mathsf T}G_\alpha\psi_\alpha(e),\\
\rho_\alpha&\ge0,
\qquad\sum_\alpha\rho_\alpha=1,\\
S_d(r,e)
&=\sum_\alpha\rho_\alpha^{(d)}K_\alpha(r,e),\\
\Delta_\alpha
&=\sum_i\eta_i\phi_\alpha(r_i)\psi_\alpha(e_i)^{\mathsf T},
\qquad
\norm{\Delta_\alpha}_F\le\varepsilon_\alpha.
\end{align}

第一式定义局部交互，第二式处理可用性和局部贡献，第三式描述当前方向化检索读出，第四式描述新实验的有界更新。重叠一致性把多个局部图约束在同一个任务层催化相容性上。

目前仍有三项重要边界。工作催化相容性尚未显式建模完整实验条件；配对级权重还没有替换当前方向化门控；有限秩更新已经进入数学内核，上层实验解析与跨图自动更新仍需独立验证。这些边界给后续研究留下了明确问题，也避免把当前模型的能力扩大到尚未验证的范围。

\end{document}
